% Options for packages loaded elsewhere
\PassOptionsToPackage{unicode}{hyperref}
\PassOptionsToPackage{hyphens}{url}
\PassOptionsToPackage{space}{xeCJK}
%
\documentclass[
  ignorenonframetext,
  aspectratio=169,
]{beamer}
\title{代謝が制約する口腔バイオフィルムの群集動態}
\subtitle{生態 ODE 推定・代謝符号 prior・空間拡張 --- 修士論文公聴会}
\author{西岡佳祐 --- NIFE}
\date{2026-11}

\usepackage{pgfpages}
\setbeamertemplate{caption}[numbered]
\setbeamertemplate{caption label separator}{: }
\setbeamercolor{caption name}{fg=normal text.fg}
\beamertemplatenavigationsymbolsempty
% Prevent slide breaks in the middle of a paragraph
\widowpenalties 1 10000
\raggedbottom
\setbeamertemplate{part page}{
  \centering
  \begin{beamercolorbox}[sep=16pt,center]{part title}
    \usebeamerfont{part title}\insertpart\par
  \end{beamercolorbox}
}
\setbeamertemplate{section page}{
  \centering
  \begin{beamercolorbox}[sep=12pt,center]{part title}
    \usebeamerfont{section title}\insertsection\par
  \end{beamercolorbox}
}
\setbeamertemplate{subsection page}{
  \centering
  \begin{beamercolorbox}[sep=8pt,center]{part title}
    \usebeamerfont{subsection title}\insertsubsection\par
  \end{beamercolorbox}
}
\AtBeginPart{
  \frame{\partpage}
}
\AtBeginSection{
  \ifbibliography
  \else
    \frame{\sectionpage}
  \fi
}
\AtBeginSubsection{
  \frame{\subsectionpage}
}
\usepackage{amsmath,amssymb}
\usepackage{lmodern}
\usepackage{iftex}
\ifPDFTeX
  \usepackage[T1]{fontenc}
  \usepackage[utf8]{inputenc}
  \usepackage{textcomp} % provide euro and other symbols
\else % if luatex or xetex
  \usepackage{unicode-math}
  \defaultfontfeatures{Scale=MatchLowercase}
  \defaultfontfeatures[\rmfamily]{Ligatures=TeX,Scale=1}
  \setmainfont[]{Noto Serif CJK JP}
  \setmonofont[]{Noto Sans Mono CJK JP}
  \ifXeTeX
    \usepackage{xeCJK}
    \setCJKmainfont[]{Noto Serif CJK JP}
  \fi
  \ifLuaTeX
    \usepackage[]{luatexja-fontspec}
    \setmainjfont[]{Noto Serif CJK JP}
  \fi
\fi
\usetheme[]{Madrid}
\usecolortheme{whale}
\usefonttheme{serif} % use mainfont rather than sansfont for slide text
% Use upquote if available, for straight quotes in verbatim environments
\IfFileExists{upquote.sty}{\usepackage{upquote}}{}
\IfFileExists{microtype.sty}{% use microtype if available
  \usepackage[]{microtype}
  \UseMicrotypeSet[protrusion]{basicmath} % disable protrusion for tt fonts
}{}
\makeatletter
\@ifundefined{KOMAClassName}{% if non-KOMA class
  \IfFileExists{parskip.sty}{%
    \usepackage{parskip}
  }{% else
    \setlength{\parindent}{0pt}
    \setlength{\parskip}{6pt plus 2pt minus 1pt}}
}{% if KOMA class
  \KOMAoptions{parskip=half}}
\makeatother
\usepackage{xcolor}
\IfFileExists{xurl.sty}{\usepackage{xurl}}{} % add URL line breaks if available
\IfFileExists{bookmark.sty}{\usepackage{bookmark}}{\usepackage{hyperref}}
\hypersetup{
  pdftitle={代謝が制約する口腔バイオフィルムの群集動態},
  pdfauthor={西岡佳祐 --- NIFE},
  hidelinks,
  pdfcreator={LaTeX via pandoc}}
\urlstyle{same} % disable monospaced font for URLs
\newif\ifbibliography
\usepackage{graphicx}
\makeatletter
\def\maxwidth{\ifdim\Gin@nat@width>\linewidth\linewidth\else\Gin@nat@width\fi}
\def\maxheight{\ifdim\Gin@nat@height>\textheight\textheight\else\Gin@nat@height\fi}
\makeatother
% Scale images if necessary, so that they will not overflow the page
% margins by default, and it is still possible to overwrite the defaults
% using explicit options in \includegraphics[width, height, ...]{}
\setkeys{Gin}{width=\maxwidth,height=\maxheight,keepaspectratio}
% Set default figure placement to htbp
\makeatletter
\def\fps@figure{htbp}
\makeatother
\setlength{\emergencystretch}{3em} % prevent overfull lines
\providecommand{\tightlist}{%
  \setlength{\itemsep}{0pt}\setlength{\parskip}{0pt}}
\setcounter{secnumdepth}{-\maxdimen} % remove section numbering
\usepackage{amsmath,amssymb}
\newcommand{\sgn}{\operatorname{sgn}}
\newcommand{\relu}[1]{\left[#1\right]_{+}}
\ifLuaTeX
  \usepackage{selnolig}  % disable illegal ligatures
\fi

\begin{document}
\frame{\titlepage}

\begin{frame}{研究の目的と概要}
\protect\hypertarget{ux7814ux7a76ux306eux76eeux7684ux3068ux6982ux8981}{}
ペリインプラント炎に関わる\textbf{口腔バイオフィルム群集の
dysbiosis}を、 群集組成の時間発展としてモデル化する（SIIRI
コンソーシアム）。

\begin{itemize}
\tightlist
\item
  臨床問題：チタンインプラント周囲で健常（commensal）から疾患（dysbiotic）へ
  遷移する微生物群集。
\item
  研究の問い：\textbf{何がこの遷移を駆動するのか} ---
  ゲノムから計算した代謝は 生態的相互作用について何を示すか。
\item
  中心的命題：\textbf{代謝が生態的相互作用の符号を制約する}。
\end{itemize}

\vspace{0.4em}

本公聴会は、生態 ODE 推定・代謝符号
prior・空間拡張を一貫した研究の弧として提示する。
\end{frame}

\begin{frame}{概念図（全体像）}
\protect\hypertarget{ux6982ux5ff5ux56f3ux5168ux4f53ux50cf}{}
\includegraphics[width=\textwidth,height=0.78\textheight]{results/figures/concept_overview_pub.png}

代謝が相互作用の\textbf{符号}を決め、dysbiosis
は\textbf{リワイヤリング＋空間再編} （Pg
の中心化・深部沈降）として現れる。
\end{frame}

\begin{frame}{研究の弧（narrative）}
\protect\hypertarget{ux7814ux7a76ux306eux5f27narrative}{}
本研究は 2 段で構成される：

\begin{enumerate}
\tightlist
\item
  \textbf{オリジナル（修士前半）} --- 5 種 in-vitro 系の \textbf{GPU
  ベイズ ODE 推定}。 TMCMC（\(10^4\) 粒子）で 4 アトラクター \textbf{CS
  / CH / DS / DH} （commensal/dysbiotic \(\times\) static/HOBIC）の
  posterior を得た。
\item
  \textbf{拡張（修士後半）} --- その枠組みを実データへ： \textbf{代謝
  AGORA 符号 prior} + \textbf{in-vivo 縦断 16S} + \textbf{空間 PDE} +
  \textbf{FISH}。
\end{enumerate}

\vspace{0.4em}

前半が力学系の骨格（アトラクター・pH 予測）を確立し、後半が
代謝の制約と空間構造でそれを実データに接地する。
\end{frame}

\begin{frame}{コアモデル（簡潔版）}
\protect\hypertarget{ux30b3ux30a2ux30e2ux30c7ux30ebux7c21ux6f54ux7248}{}
絶対量 \(x_i\) の一般化 Lotka--Volterra：
\[\dot{x}_i = x_i\Big(b_i + \sum_{j=1}^{S} A_{ij}\,x_j\Big),\qquad S=10 .\]

単体 \(\sum_i\varphi_i=1\) 上の replicator / Hamilton 形（16S
は組成データ）：
\[\dot{\varphi}_i = \varphi_i\Big[(A\varphi+b)_i - \varphi^{\!\top}(A\varphi+b)\Big].\]

\textbf{符号 prior}：AGORA cross-feeding フラックス \(F_{ij}\) から
\[P_{ij}=\operatorname{sgn}(F_{ij})\in\{-1,0,+1\}\] を作り、\(A\)
の\textbf{符号のみ}を制約（大きさは設計上捨てる）。
\end{frame}

\begin{frame}[fragile]{データフロー（実データ）}
\protect\hypertarget{ux30c7ux30fcux30bfux30d5ux30edux30fcux5b9fux30c7ux30fcux30bf}{}
\includegraphics[width=\textwidth,height=0.66\textheight]{results/figures/pipeline_overview_pub.png}

実 Dieckow データ：guild \(\varphi\)（10 患者 \(\times\) 3 週）→ AGORA
符号 prior → gLV/Hamilton フィット → LOO-CV 検証 → 4 アトラクター。
\texttt{fit\_*.json} が交換フォーマット。
\end{frame}

\begin{frame}{AGORA 符号 prior パイプライン}
\protect\hypertarget{agora-ux7b26ux53f7-prior-ux30d1ux30a4ux30d7ux30e9ux30a4ux30f3}{}
\includegraphics[width=\textwidth,height=0.64\textheight]{results/fig2_agora_pipeline.png}

\begin{enumerate}
[(A)]
\tightlist
\item
  手順；(B) 層追加で制約集合 \(|\mathcal{M}|\) が \(10\to22\to58\)
  （AGORA L3 で +36）；(C) 出来上がる符号 prior 行列
  \(P=\operatorname{sgn}(F)\)。
\end{enumerate}
\end{frame}

\begin{frame}{経験的 sign agreement（naive 推定）}
\protect\hypertarget{ux7d4cux9a13ux7684-sign-agreementnaive-ux63a8ux5b9a}{}
\includegraphics[width=\textwidth,height=0.55\textheight]{results/fig3_agora_sign_validation.png}

\[\mathrm{SA}=\frac{1}{|\mathcal{M}|}\sum_{(i,j)\in\mathcal{M}}\mathbb{1}\!\left[\operatorname{sgn}(\hat A_{ij})=P_{ij}\right]=92\%.\]
\textcolor{red}{これは過大評価 — prior は符号縮退している（次スライドで誠実に検証）。}
\end{frame}

\begin{frame}{誠実な検証}
\protect\hypertarget{ux8aa0ux5b9fux306aux691cux8a3c}{}
naive 92\% は prior の符号縮退による過大評価。正しい対照は permutation：

\begin{itemize}
\tightlist
\item
  \textbf{cross-feeding 方向のみ}が独立に検証される：
  \[p=4\times10^{-4}\quad(z=+3.79,\ n_\pi=10^4).\]
\item
  \textbf{競争方向は検証されない}（chance レベル）。
\item
  \textbf{2 コホート}（Dieckow \(\times\) Duran-Pinedo）を prior 抜きで比べ、
  強ペア符号が \textbf{89\%} 一致（\(p\approx0.02\)）。
\end{itemize}

\vspace{0.3em}

\textcolor{red}{AGORA prior は 16S 力学で再現されない $\Rightarrow$ モデリング選択であり、
データで確認済みではない。}
\end{frame}

\begin{frame}{ネットワークの視点}
\protect\hypertarget{ux30cdux30c3ux30c8ux30efux30fcux30afux306eux8996ux70b9}{}
推定 \(A\) の構造解析（class レベル）：

\begin{itemize}
\tightlist
\item
  \textbf{Veillonella（Negativicutes）が代謝のシンク} ---
  乳酸などを集約する受け手。
\item
  古典的な \textbf{Pg-keystone / Fn-bridge は class
  レベルでは支持されない}。
\item
  dysbiosis で \textbf{Pg が中心化}する（固有中心性 \(0.32\to0.51\)）。
\item
  同時に \textbf{S. oralis--Veillonella の相利共生が崩れる}。
\end{itemize}

\vspace{0.4em}

dysbiosis
は単なる組成変化でなく、相互作用ネットワークの\textbf{リワイヤリング}。
\end{frame}

\begin{frame}{空間拡張 --- FISH 深さプロファイル（CH vs DH）}
\protect\hypertarget{ux7a7aux9593ux62e1ux5f35-fish-ux6df1ux3055ux30d7ux30edux30d5ux30a1ux30a4ux30ebch-vs-dh}{}
\includegraphics[width=\textwidth,height=0.5\textheight]{results/diffusion_fit/zprofiles_all_ti_overlay.png}

CLSM-FISH から得た\textbf{深さ分解 5
種プロファイル}。commensal-HOBIC（CH）と
dysbiotic-HOBIC（DH）で深さ方向の構造が異なる。
\end{frame}

\begin{frame}{dysbiosis = 空間的再編成}
\protect\hypertarget{dysbiosis-ux7a7aux9593ux7684ux518dux7de8ux6210}{}
\includegraphics[width=\textwidth,height=0.52\textheight]{results/diffusion_fit/depth_niche.png}

dysbiosis は組成変化に留まらず\textbf{空間的再編成}： \emph{P.
gingivalis} が深部へ沈む（6 日目から \textbf{+30 µm}、嫌気ニッチ）。
\end{frame}

\begin{frame}{新規：空間的 cross-feeding テスト}
\protect\hypertarget{ux65b0ux898fux7a7aux9593ux7684-cross-feeding-ux30c6ux30b9ux30c8}{}
\includegraphics[width=\textwidth,height=0.48\textheight]{results/diffusion_fit/spatial_crossfeeding.png}

乳酸ペア \textbf{S. oralis（生産者）\(\to\)
Veillonella（消費者）}が空間的に層状化： Veillonella
が\textbf{より浅い}（10/10 サンプル、Wilcoxon \(p=0.002\)）。
生産者が基底、消費者がその上 --- \textbf{上向きの乳酸拡散}と整合。
\end{frame}

\begin{frame}{反応拡散 PDE}
\protect\hypertarget{ux53cdux5fdcux62e1ux6563-pde}{}
0D gLV を深さ \(z\) 方向の反応拡散へ拡張：
\[\partial_t\varphi_i = D_i\,\partial_z^2\varphi_i - u\,\partial_z\varphi_i + R_i(\varphi),\]
\(R_i\) は replicator 反応項（\(A,b\) は時間 fit から固定）、
\(D_i\)（拡散）と \(u\)（移流）のみを深さプロファイルから推定。

\vspace{0.4em}

\textcolor{red}{現状は preliminary — HPC でのパラメータ sweep が収束中。}
空間 fit は時間動態と独立な観測量（FISH 深さ）でモデルを接地する。
\end{frame}

\begin{frame}{pH 予測（オリジナル 5 種研究・独立検証）}
\protect\hypertarget{ph-ux4e88ux6e2cux30aaux30eaux30b8ux30caux30eb-5-ux7a2eux7814ux7a76ux72ecux7acbux691cux8a3c}{}
\includegraphics[width=\textwidth,height=0.52\textheight]{results/figures/fig_ph_validation.png}

5 種アトラクター posterior から pH を前向きに予測： 独立
\(R^2=0.78\)、RMSE \(0.13\)、LOO \(R^2=0.92\)。 \textbf{pH
はキャリブレーションに使っていない} --- 真に独立な検証。
\end{frame}

\begin{frame}{結論}
\protect\hypertarget{ux7d50ux8ad6}{}
\begin{enumerate}
\tightlist
\item
  \textbf{代謝が生態的相互作用の符号を制約する}
  （\(P_{ij}=\operatorname{sgn}(F_{ij})\)；大きさは使わない）。
\item
  dysbiosis は\textbf{リワイヤリング + 空間的再編成}であり、 \emph{P.
  gingivalis} の深部沈降（+30 µm）を伴う。
\item
  その主張は \textbf{2 コホート}（Dieckow \(\times\) Duran-Pinedo, 89\%,
  \(p\approx0.02\)）と \textbf{独立な機械論経路}（COMETS
  dFBA）で再現される。
\item
  オリジナルの 5 種 ODE は pH を独立に予測（\(R^2=0.78\), LOO
  \(R^2=0.92\)）。
\end{enumerate}
\end{frame}

\begin{frame}{今後（Outlook）}
\protect\hypertarget{ux4ecaux5f8coutlook}{}
\begin{itemize}
\tightlist
\item
  \textbf{空間的 out-of-sample 検証} ---
  別サンプルで深さプロファイルを予測。
\item
  \textbf{2D / 3D PDE} --- \(z\) 方向だけでなく面内・立体構造へ。
\item
  \textbf{PINN 逆問題} --- 物理情報ニューラルネットで \(D_i,u\)
  を直接推定。
\item
  \textbf{GDI / Joshi 臨床検証} --- 臨床メタデータ統合（dysbiosis
  index）。
\end{itemize}

\vspace{0.4em}

代謝が符号を制約し、dysbiosis がリワイヤリング＋空間再編であるという
描像を、空間検証と臨床指標によってさらに検証する。
\end{frame}

\begin{frame}[fragile]{まとめ（Take-home）}
\protect\hypertarget{ux307eux3068ux3081take-home}{}
\begin{center}
\large
代謝は生態的相互作用の\textbf{符号}を制約する。\\[0.3em]
dysbiosis は\textbf{リワイヤリング＋空間再編}である。\\[0.3em]
\textbf{2 コホート}＋\textbf{機械論的 dFBA}＋\textbf{独立な pH 予測}で再現。
\end{center}

\vspace{0.6em}

すべてを 1 つの \textbf{10 ギルド分類}と \textbf{\texttt{fit\_*.json}}
が束ねる。
\end{frame}

\end{document}
